\begin{abstract}
Transformer scaling laws (\emph{the empirical curves that predict how a language model's performance grows with its parameter count and the amount of text it is trained on}) are seemingly fitted on analytic-leaning languages such as English and Mandarin, whose words carry little grammatical marking on their written surface. This could genuinely be due to the fact that the AI market is dominated by China and Anglophone regions of the world at the time of writing. This paper asks whether languages at the opposite end of the typological spectrum, whose words inflect richly and write that inflection out in full, can turn their grammatical structure into a measurable computational efficiency factor.

The mechanism we study is \emph{grammar-aware tokenisation}. Instead of chopping text into frequent character fragments by statistics alone (\emph{byte-pair encoding, or BPE for short, which never consults grammar}), we hand the model a grammatical analysis of each word: its part of speech and the bundle of inflectional features it carries (\emph{tense, person, gender, number, case, and so on}). This analysis enters the network as a feature embedding summed with the ordinary token embedding at the input layer. Arabic's templatic nature means that part of that contextual information is stored deeper than the surface, and accordingly it has two further streams exposing the consonantal root and the \emph{wazn} pattern, more on that later. The hypothesis, designed to be tested rather than assumed, is that exposing the morphological structure of our test languages rebates part of their transformers' parameter budgets, and that the rebate scales with how much grammar the writing system already makes visible: a language whose grammar indicates lower morphological information density should gain little; one that indicates more should gain more.

We try to make this precise through a single quantity, $\rho(L) \cdot H(L)$, where $H(L)$ is the Shannon entropy of the grammatical features a word form carries (\emph{a bit-valued measure of how much grammar one word packs}) and $\rho(L)$ is the share of that information legible from the written word itself, without lexical lookup. A rule-based grammar engine computes both before any model is trained, so the cross-linguistic ordering is registered in advance: $\rho \cdot H$ ranks the four languages Mandarin $<$ English $<$ Turkish $<$ Arabic, with values $0.16$, $1.07$, $2.98$, and $3.47$ on the four-language corpora. The same ordering is already visible in the raw inventories the engine produces: across the three inflecting languages the number of distinct grammatical feature bundles attested follows the same rank (\emph{roughly $21$ for English, $469$ for Turkish, and over $1{,}400$ for Arabic}), so the ranking is legible before entropy or surface-recoverability is computed.

The study trains four matched pairs of transformers, a baseline and a grammar-aware model for each of Mandarin, English, Turkish, and Arabic, four languages chosen to span the typological spectrum from isolating through analytic and agglutinative to templatic. Each model holds roughly $3.2 \times 10^7$ parameters in the transformer the two regimes share and is trained on about $6.4 \times 10^8$ tokens (\emph{close to twenty per parameter, the compute-optimal ratio}) drawn from some one billion tokens of mixed encyclopaedic and web-crawl text, five million sentences per language. Because the baseline and grammar-aware tokenisers cut the same text into different numbers of pieces, a per-token loss is not comparable between them; we therefore report bits per character, which scores the cost of the text itself.

Mandarin, the isolating control, behaves as predicted: its baseline and grammar-aware models finish within $0.004$ bits per character of each other, the near-null a language that writes almost no grammar on its surface should show. English, the analytic case, shows no rebate either at this scale, its grammar-aware model finishing $0.036$ bits per character \emph{worse} than baseline. Turkish, the agglutinative case where the hypothesis predicts a clear rebate, also shows none at this scale: its grammar-aware model finishes $0.017$ bits per character worse than baseline. Arabic, the templatic case, is the exception: its grammar-aware model finishes $0.0192 \pm 0.0023$ bits per character \emph{better} than baseline (\emph{mean over four seeds}), the only positive rebate of the four. The pattern that emerges is narrower than the framework's $\rho \cdot H$ ordering. The rebate appears only for the language whose grammar is least legible from the written surface (Arabic, the lowest $\rho$, whose root-and-pattern system hides structure off the surface), and is absent wherever a compute-optimal model can read the grammar off the text for itself (Mandarin, English, Turkish). One caveat could temper the Arabic result: its grammar-aware model also carries the most additional embedding parameters (the root and \emph{wazn} streams raise its parameter count by roughly half). A parameter-matched ablation settles it: scrambling the content of those two streams while holding their size fixed makes the model finish \emph{below} its own baseline, and the real streams beat the scrambled ones by $0.043$ bits per character, so the gain is attributable to templatic structure rather than to added capacity.

A smaller three-language pilot first surfaced an ordering by $\rho \cdot H$ and motivated this study; the four-language test overturns it. At compute-optimal, transformer-matched scale (\emph{the grammar-aware models carry more total parameters, which only makes the null results more conservative}), the rebate vanishes for the three languages whose grammar is surface-legible, Turkish included, the agglutinative case the pilot most favoured, because a model with enough data learns their grammar from the text itself, making an explicit morphological analysis redundant. It survives only for Arabic, whose templatic morphology hides grammar off the surface, and a shuffle-control ablation indicates that survival is structural rather than a parameter artefact. Where the lever exists at all, it appears to track the grammar a language hides off the surface rather than the recoverable $\rho \cdot H$ the framework first proposed; with a single positive language, this is a conjecture for a larger study, not a law we establish. The practical promise, on-device modelling of grammar-rich languages at parameter counts English does not reach (\emph{a frontier that compact edge models such as Google's Gemma~4 E4B, a 2026 on-device model that appeared while this research was underway and uses per-layer embeddings to reach roughly four billion effective parameters while still fitting on a phone, are already approaching}), therefore narrows to the templatic case, where it is strongest and where current pipelines serve worst.
\end{abstract}
